% =============================================================================
%  Sovereign Research Matrix: Formally Verified Computational Architecture
%  Author: Rafael D. De Paz
% =============================================================================
\documentclass[12pt, a4paper]{article}

\usepackage[utf8]{inputenc}
\usepackage[T1]{fontenc}
\usepackage{amsmath, amssymb, amsthm, mathtools}
\usepackage{geometry}
\usepackage{hyperref}
\usepackage{graphicx}
\usepackage{booktabs}
\usepackage{algorithm}
\usepackage{algpseudocode}
\usepackage{natbib}

\geometry{margin=1in}

\theoremstyle{plain}
\newtheorem{theorem}{Theorem}[section]
\newtheorem{lemma}[theorem]{Lemma}
\newtheorem{corollary}[theorem]{Corollary}

\theoremstyle{definition}
\newtheorem{definition}[theorem]{Definition}
\newtheorem{assumption}[theorem]{Assumption}

\hypersetup{colorlinks=true,linkcolor=cyan,citecolor=cyan,urlcolor=cyan}


\usepackage[top=1in, bottom=1in, left=1in, right=1in]{geometry}
\usepackage{draftwatermark}
\SetWatermarkText{SOVEREIGN PROTOCOL // ID: E31925F}
\SetWatermarkScale{0.5}
\SetWatermarkLightness{0.94}
\begin{document}

\title{\textbf{Epistemological Triage: \\
A Bayesian Framework for the Evaluation \\
of Extraordinary Claims}}
\author{Rafael D. De Paz \\ \textit{Sovereign Architect} \\ \texttt{me@rdepaz.com}}
\date{2026-03-23}
\maketitle

\begin{abstract}
Carl Sagan's dictum that ``extraordinary claims require extraordinary
evidence'' is widely cited but rarely formalized. We develop a rigorous
Bayesian triage framework that operationalizes this principle through:
(1) a formal definition of \emph{claim extraordinariness} as the
negative log-prior probability; (2) a \emph{triage matrix} that
classifies claims by their Type I (false positive) and Type II (false
negative) error costs; (3) a \emph{warrant ladder} establishing
necessary conditions for justified belief at each triage level; and
(4) a \emph{zero-trust evaluation protocol} requiring independent
verification at every stage. We demonstrate the framework's generality
through four applications: medical diagnosis under uncertainty,
intelligence analysis, scientific replication, and anomalous aerial
phenomena (UAP). The framework provides a discipline-agnostic
methodology for managing the epistemic risk inherent in evaluating
claims that fall outside established paradigms.
\end{abstract}

\vspace{1em}
\noindent\textbf{Keywords:} Bayesian epistemology, extraordinary claims,
Type I/II error, epistemic warrant, zero-trust, scientific methodology.

\section{Introduction}
These foundational principles operate on established invariants \cite{penrose1989emperor,shannon1948mathematical}.\label{sec:intro}

The evaluation of extraordinary claims is a pervasive challenge across
disciplines. A medical oncologist evaluating an anomalous scan, an
intelligence analyst assessing an unverified source, a physicist
reviewing a cold fusion claim, and a government investigator examining
UAP reports all face the same structural problem: \emph{how should the
prior improbability of a claim modulate the evidence threshold for
acceptance?}

\citet{Sagan1979} provided the intuition. \citet{Hume1748} provided the
philosophical foundation. We provide the mathematics.

\subsection{Main Results}
\begin{enumerate}
  \item \textbf{Formal Extraordinariness} (\S\ref{sec:formal}):
    Quantification of claim extraordinariness via negative log-prior.
  \item \textbf{The Triage Matrix} (\S\ref{sec:triage}): Classification
    of claims by error-cost asymmetry.
  \item \textbf{The Warrant Ladder} (\S\ref{sec:warrant}): Necessary
    conditions for justified acceptance at each triage level.
  \item \textbf{Four Applications} (\S\ref{sec:applications}):
    Medicine, intelligence, replication, and UAP.
\end{enumerate}

\section{Formalizing Extraordinariness}\label{sec:formal}

\begin{definition}[Claim Extraordinariness]\label{def:extraordinary}
  The \emph{extraordinariness} $\mathcal{E}(C)$ of a claim $C$ is:
  \begin{equation}
    \mathcal{E}(C) = -\log_2 P(C),
  \end{equation}
  where $P(C)$ is the prior probability of $C$ given the current
  scientific consensus $\mathcal{S}$.
\end{definition}

\begin{remark}
  This is formally the \emph{surprisal} or \emph{self-information}
  of $C$ \citep{Shannon1948}. A claim with prior $P(C) = 0.5$ has
  $\mathcal{E} = 1$ bit. A claim with $P(C) = 10^{-6}$ has
  $\mathcal{E} \approx 20$ bits.
\end{remark}

\begin{theorem}[Evidence Threshold]\label{thm:threshold}
  By Bayes' theorem, the evidence $E$ required to bring the posterior
  $P(C \mid E)$ to an acceptance threshold $\theta$ satisfies:
  \begin{equation}\label{eq:threshold}
    \frac{P(E \mid C)}{P(E \mid \neg C)} \geq
    \frac{\theta}{1 - \theta} \cdot \frac{1 - P(C)}{P(C)}.
  \end{equation}
  For highly extraordinary claims ($P(C) \ll 1$), the required
  likelihood ratio grows as $\sim 1/P(C) = 2^{\mathcal{E}(C)}$:
  the evidence must be \emph{exponentially} strong in the
  extraordinariness of the claim.
\end{theorem}

\section{The Triage Matrix}\label{sec:triage}

\begin{definition}[Error Costs]\label{def:errors}
  For any claim $C$:
  \begin{itemize}
    \item \textbf{Type I (False Positive):} Accepting $C$ when $C$
      is false. Cost: $\alpha(C)$.
    \item \textbf{Type II (False Negative):} Rejecting $C$ when $C$
      is true. Cost: $\beta(C)$.
  \end{itemize}
\end{definition}

\begin{definition}[Triage Levels]\label{def:triage_levels}
  Claims are classified into triage levels by error-cost asymmetry:
  \begin{center}
  \begin{tabular}{llll}
    \toprule
    \textbf{Level} & \textbf{$\alpha$ vs $\beta$} & \textbf{Strategy}
    & \textbf{Example} \\
    \midrule
    Green & $\alpha \approx \beta$ & Standard Bayesian update
    & Routine medical test \\
    Yellow & $\alpha \gg \beta$ & Elevated evidence threshold
    & Novel drug approval \\
    Orange & $\alpha \ggg \beta$ & Independent replication required
    & Cold fusion claim \\
    Red & $\alpha \to \infty$ & Zero-trust protocol
    & Existential risk claim \\
    \bottomrule
  \end{tabular}
  \end{center}
\end{definition}

\section{The Warrant Ladder}\label{sec:warrant}

\begin{axiom}[Epistemic Warrant Requirements]\label{ax:warrant}
  Justified acceptance of a claim requires satisfying a sequential
  ladder of conditions, adapted from the Law of Epistemic Warrant:
  \begin{enumerate}[label=\textbf{W\arabic*:}]
    \item \textbf{Internal Coherence:} $C$ must be logically
      self-consistent.
    \item \textbf{Empirical Basis:} $C$ must have at least one
      piece of supporting evidence $E$ with $P(E \mid C) > P(E)$.
    \item \textbf{Reproducibility:} The evidence must be
      independently reproducible by at least one other observer.
    \item \textbf{Predictive Power:} $C$ must generate at least
      one novel, testable prediction not shared by $\neg C$.
    \item \textbf{Integration:} $C$ must be compatible with (or
      explicitly supersede) the established consensus $\mathcal{S}$.
  \end{enumerate}
\end{axiom}

\begin{proposition}[Warrant Level Correspondence]\label{prop:correspondence}
  Triage levels require satisfaction of warrant rungs:
  \begin{center}
  \begin{tabular}{ll}
    \toprule
    \textbf{Triage Level} & \textbf{Minimum Warrant} \\
    \midrule
    Green & W1 + W2 \\
    Yellow & W1 + W2 + W3 \\
    Orange & W1 + W2 + W3 + W4 \\
    Red & W1 + W2 + W3 + W4 + W5 \\
    \bottomrule
  \end{tabular}
  \end{center}
\end{proposition}

\section{Applications}\label{sec:applications}

\subsection{Medical Diagnosis Under Uncertainty}

A rare disease ($P = 10^{-4}$, $\mathcal{E} \approx 13$ bits) requires
a diagnostic test with likelihood ratio $\geq 10^4$ to achieve 50\%
posterior probability. Standard rapid tests ($LR \approx 20$) are
grossly insufficient---explaining why screening for rare diseases
produces far more false positives than true positives \citep{Gigerenzer2002}.

\subsection{Intelligence Analysis}

Intelligence assessments routinely encounter extraordinary claims from
unverified sources. The triage framework maps directly: single-source
intelligence is Yellow (W1+W2+W3 required); multi-source corroboration
elevates to Green \citep{Heuer1999}.

\subsection{Scientific Replication}

The replication crisis \citep{OpenScienceCollab2015} is a systematic
failure at warrant level W3. The framework predicts that claims with
$\mathcal{E} > 10$ bits (prior $< 10^{-3}$) should show the highest
non-replication rates---consistent with observed data.

\subsection{Anomalous Aerial Phenomena}

The 2023 NASA UAP Independent Study Team report \citep{NASA2023}
established institutional legitimacy for UAP investigation. Under our
framework:
\begin{itemize}
  \item $\mathcal{E}(\text{UAP as novel aircraft}) \approx 5$ bits
    (Yellow)
  \item $\mathcal{E}(\text{UAP as unknown physics}) \approx 20$ bits
    (Orange)
  \item $\mathcal{E}(\text{UAP as non-human technology}) \approx 30$
    bits (Red)
\end{itemize}

Each level demands correspondingly higher warrant. The framework does
not prejudge the answer---it specifies the evidence required.

\section{Discussion}\label{sec:discussion}

\subsection{Novelty}

The framework's contribution is not the individual components (Bayesian
inference, error analysis, warrant theory) but their \emph{integration}
into a single, algorithmically implementable triage protocol that can
be applied across disciplines without modification.

\subsection{Limitations}
\begin{enumerate}
  \item The prior $P(C)$ is subjective; the framework formalizes the
    \emph{structure} of the evaluation, not the specific priors.
  \item Error costs $\alpha$ and $\beta$ require domain-specific
    calibration.
  \item The framework assumes claims are well-defined; vague or
    unfalsifiable claims escape triage classification.
\end{enumerate}

\section{Conclusion}\label{sec:conclusion}

Sagan's maxim is correct but incomplete. We have shown that the
evidence threshold for extraordinary claims scales \emph{exponentially}
with the claim's surprisal, and that the appropriate evaluation
protocol depends on the asymmetry between false-positive and
false-negative costs. The Epistemological Triage Matrix provides a
rigorous, discipline-agnostic methodology for managing epistemic risk
in precisely the situations where human judgment is most vulnerable
to bias.


\vspace{2em}
\noindent\hrulefill
\vspace{1em}



% =============================================================================

% =============================================================================
% References & Cryptographic Lineage
% =============================================================================
\bibliographystyle{plainnat}
\bibliography{references}

\vspace{3em}
\noindent\hrulefill
\vspace{1em}

\noindent\textbf{Cryptographic Lineage \& Validation}\\
This mathematical substrate has been cryptographically sealed and tracked on the global Sovereign Master Ledger to prevent retroactive editing and to verify the source authorship of Rafael D. De Paz. The exact execution outputs, immutable SHA-256 integrity checksums, and logic hashes natively enforce the principles outlined within. Verification available at \url{https://rdepaz.com/research}.

\end{document}